Die Wärmepumpe nimmt Wärme bei der Aussentemperatur $T_\text{a} = \Ta$ auf und gibt sie bei der Innentemperatur $T_\text{i} = \Ti$ ab. Ihre maximale Leistungsziffer ist (FoTaBe S.~172)
\begin{align}
 \eta_\text{WP} &= \ewpF \\
   &= \frac{\Ti}{\Ti-\Ta} \\
   &= \ewp \approx \result{\ewpP{0}{2}}
\end{align}
Eine ideale Wärmepumpe liefert also rund achtmal so viel Wärme ins Haus, wie sie an elektrischer Energie aufnimmt.
